% Section 3 — A commitment protocol for one proposal
\section{A commitment protocol for one move}
\label{sec:protocol}

\begin{definition}{The Zero contract}
\label{def:zero}
A proposal satisfies the Zero contract when, for the problem at hand,
\begin{itemize}
\item $K=1$: exactly one complete candidate is committed---no sampling several and keeping one, no
beam over competing candidates, no best-of-$B$;
\item $E=0$: no candidate is executed, compiled for feedback, simulated or scored, by a verifier or by
a learned critic, before commitment;
\item $R=0$: nothing is revised after any feedback; a repaired candidate is a new attempt and the
first attempt's failure stands.
\end{itemize}
Permitted: frozen normalization, declared tokenization, direct feature computation, fixed typed
grammar masks, one sampled or greedy decoding trajectory with every step counted, and deterministic
serialization. Counters are recorded by an isolated supervisor and never taken from the proposer's
own report.
\end{definition}

\subsection{Four tracks}

The contract separates four settings that are routinely blurred (\cref{tab:tracks,fig:protocol}).
\textbf{C0} is the cold-start artifact: the player sees the specification and authorized
observations---for a learning procedure, the training rows---and commits one candidate with no new
probe of the target task. \textbf{CH} is the history-conditioned decision: every compared method
receives the same frozen, disclosed history of earlier evaluations and emits one new decision; the
history is charged to the workflow that produced it. \textbf{SEQ} is a sequential controller: one
proposal per step, with new observations in between, reported as $T$ decisions and their total
evaluation budget, never as ``one inference''. \textbf{S} is search-assisted: multiple candidates,
ranking, evaluation or repair at a declared budget $B$. A learned critic may be used in training or in
an explicitly labelled S arm; its deployment use, not its existence, decides the track.

\begin{table}[t]
\caption{Four tracks. Every result is labelled with its track, and no comparison crosses tracks
without saying so.}
\label{tab:tracks}
\centering\tablesf
\begin{tabularx}{\linewidth}{@{}l>{\raggedright\arraybackslash}p{0.27\linewidth}LL@{}}
\toprule
Track & Starting information & Permitted before commitment & Correct description \\
\midrule
C0 cold start & specification and authorized observations & one trajectory; no probes, execution, feedback or revision & first artifact from the supplied state \\
CH history & a frozen, disclosed history, identical across methods & one new proposal; no look-ahead evaluation & one decision given an existing history \\
SEQ sequential & an initial state, then accumulating outcomes & one proposal per step & $T$ decisions and their evaluation budget \\
S search & declared state and search budget $B$ & several candidates, ranking, evaluation, repair & search-enhanced performance at budget $B$ \\
\bottomrule
\end{tabularx}
\end{table}

\begin{figure}[t]
\centering
% Figure — the information and commitment boundary, with the four tracks. Conceptual; no data.
\begin{tikzpicture}[x=1cm, y=1cm,
  st/.style={ezbox, font=\sffamily\scriptsize, text width=1.75cm, minimum height=0.85cm, inner sep=2.5pt},
  sg/.style={ezstage, font=\sffamily\scriptsize, text width=1.75cm, minimum height=0.85cm, inner sep=2.5pt},
  lane/.style={font=\sffamily\sbseries\scriptsize, text=ezaccent, anchor=west},
  note/.style={font=\sffamily\scriptsize, text=ezgrey, anchor=south west, inner sep=1.5pt},
  dot/.style={circle, fill=ezink, inner sep=0pt, minimum size=4pt},
  run/.style={ezink, line width=0.7pt, -{Stealth[length=2mm]}}]
  % backdrops
  \fill[ezaccentsoft, rounded corners=3pt] (-0.15,0.85) rectangle (7.05,-4.35);
  \fill[ezpaper, rounded corners=3pt] (7.35,0.85) rectangle (14.95,-4.35);
  \node[font=\sffamily\sbseries\scriptsize, text=ezaccent, anchor=north west] at (-0.05,0.8) {known to the player};
  \node[font=\sffamily\sbseries\scriptsize, text=ezgrey, anchor=north east] at (14.85,0.8) {verification only};
  % generic pipeline
  \node[st] (s) at (0.95,-0.4) {state $s$\\spec, data, history};
  \node[sg] (e) at (3.05,-0.4) {encode\\$z=E_{\phi,d}(s)$};
  \node[sg] (p) at (5.25,-0.4) {propose\\$h\to a$, one path};
  \node[st] (x) at (9.0,-0.4) {execute\\fit, run, build};
  \node[st] (v) at (11.1,-0.4) {verify $V_d$\\withheld data};
  \node[st] (o) at (13.3,-0.4) {outcome\\$(v,y,c)$, $U$};
  \draw[ezarrow] (s) -- (e); \draw[ezarrow] (e) -- (p); \draw[ezarrow] (p) -- (x);
  \draw[ezarrow] (x) -- (v); \draw[ezarrow] (v) -- (o);
  % commitment line
  \draw[ezwarm, line width=1pt] (7.2,1.05) -- (7.2,-4.5);
  \node[font=\sffamily\sbseries\scriptsize, text=ezwarm, fill=white, inner sep=1.5pt] at (7.2,1.05) {commit: hash, freeze};
  % lanes
  \foreach \y in {-1.55,-2.3,-3.05,-3.95} \draw[ezrule, line width=0.3pt] (0.0,\y) -- (14.8,\y);
  % C0
  \node[lane] at (0.0,-1.4) {C0};
  \node[dot] at (1.6,-1.55) {}; \draw[run] (1.6,-1.55) -- (13.0,-1.55);
  \node[note] at (1.75,-1.55) {specification and authorized observations; no new probes};
  % CH
  \node[lane] at (0.0,-2.15) {CH};
  \filldraw[fill=ezamber!25, draw=ezamber, line width=0.4pt] (0.78,-2.42) rectangle (1.42,-2.18);
  \node[dot] at (1.6,-2.3) {}; \draw[run] (1.6,-2.3) -- (13.0,-2.3);
  \node[note] at (1.75,-2.3) {frozen history $H_t$, identical for every method, charged to its workflow};
  % SEQ
  \node[lane] at (0.0,-2.9) {SEQ};
  \node[dot] at (1.6,-3.05) {}; \draw[run] (1.6,-3.05) -- (13.0,-3.05);
  \draw[ezink, line width=0.55pt, dashed, -{Stealth[length=1.8mm]}] (13.15,-3.05) -- (13.15,-3.42) -- (1.6,-3.42) -- (1.6,-3.13);
  \node[note] at (1.75,-3.05) {one move per step};
  \node[font=\sffamily\scriptsize, text=ezgrey, fill=ezpaper, inner sep=1.2pt] at (10.2,-3.42) {$s_{t+1}=s_t\oplus(a_t,v_t,y_t,c_t)$};
  % S
  \node[lane] at (0.0,-3.8) {S};
  \node[dot] at (1.6,-3.95) {}; \draw[run] (1.6,-3.95) -- (13.0,-3.95);
  \draw[ezgrey, line width=0.55pt, rounded corners=2pt, -{Stealth[length=1.5mm]}] (8.4,-3.95) -- (8.4,-3.72) -- (6.0,-3.72) -- (6.0,-3.9);
  \node[note] at (1.75,-3.95) {$B$ candidates scored};
  \node[note, anchor=south west] at (8.55,-3.95) {up to $B$ returns, each one counted};
\end{tikzpicture}

\caption{The information boundary. Everything left of the commitment line is available to the
player; the candidate is hashed and frozen before any candidate-specific execution; everything to the
right is verification. In C0 and CH the game ends at the verifier. In SEQ the verified outcome
becomes part of the next state, and each step is a new move under the same contract. In S the player
may cross the line repeatedly before committing, and every crossing is counted.}
\label{fig:protocol}
\end{figure}

\subsection{What counts as computation}

One proposal is not one forward pass. An autoregressive decoder makes many network evaluations while
writing one program; an encoder may process many chunks. All can belong to one declared proposal
transaction, but they are different computational claims, so the supervisor records them separately
(\cref{tab:counters}). A parallel batch is not free because it arrives in one request, and missing
telemetry is reported as unknown, never as zero.

\begin{table}[t]
\caption{Counters recorded for every proposal by the supervisor (\cref{app:formal} gives the
schema).}
\label{tab:counters}
\centering\tablesf
\begin{tabularx}{\linewidth}{@{}>{\raggedright\arraybackslash}p{0.36\linewidth}L@{}}
\toprule
Counter & Meaning \\
\midrule
\code{proposal\_transactions} & complete user-visible proposal requests \\
\code{candidate\_trajectories}, \code{complete\_candidates} & competing paths and complete candidates generated, including discarded ones (strict tracks: exactly 1) \\
\code{encoder\_forward\_calls}, \code{proposer\_forward\_calls}, \code{decoder\_steps} & network computation, with the batching convention \\
\code{new\_target\_probe\_evaluations} & new evaluations of the target task acquired during the request (C0: 0) \\
\code{candidate\_executions\_before\_commit}, \code{surrogate\_candidate\_scores} & compiler, simulator, runtime or critic feedback on candidates (strict tracks: 0) \\
\code{repair\_or\_retry\_count}, \code{external\_model\_or\_tool\_calls} & revisions and outside services used during proposal (strict tracks: 0) \\
\code{historical\_evaluations\_supplied} & observations and costs supplied as input history (CH and SEQ) \\
\code{postcommit\_verification\_work} & all later fitting, testing, simulation and physical evaluation \\
\bottomrule
\end{tabularx}
\end{table}

\subsection{Commit, execute, verify}

The candidate includes everything that determines its behaviour: source or construction sequence,
structural choices, constants, dependencies, seeds and the execution recipe. The supervisor binds
them in a proposal receipt before any feedback exists; the verifier then runs the unchanged
candidate, and only the verifier holds the withheld data.

\emph{Learning procedures} are fitted on the authorized training rows after commitment; ordinary
fitting is disclosed execution cost. Hyperparameter sweeps, model selection or validation-chosen
ensembles inside \code{fit} are \emph{search-in-fit} and are reported in a separate cohort; the
primary cohort forbids them. \emph{CAD programs} are built, meshed and simulated after commitment; a
build failure is the proposal's failure, and any geometry healing is a fixed declared transformation
applied identically to every method or reported as a repair-assisted result. \emph{Optimization
points} are fixed, with any fidelity choice, before the next objective value exists, and a rejected
point is not replaced for free.

\subsection{Failure, abstention and fallback}

The scientific primary keeps the first proposal's failure. A deployed system may fall back to a safe
conventional method, and we then report both the first-proposal result and the total service result,
fallback cost included. Abstention is also an output: we report coverage and utility over all
requested tasks. Invalid syntax, invalid types, non-finite values, timeouts, infeasible designs,
evaluator faults and missing measurements are distinguished; unknown outcomes are neither failures
nor successes, and infrastructure retries follow a rule fixed before any method runs.
